\documentclass[11pt]{letter}

\usepackage{geometry}
\geometry{margin=1in}
\usepackage{setspace}
\usepackage{hyperref}

\signature{Decory Edwards}
\address{Decory Edwards \\ Johns Hopkins University \\ Department of Economics}

\begin{document}

\begin{letter}{Hiring Committee \\ Susquehanna International Group}

\opening{Dear Hiring Committee,}

I am writing to apply for the Quantitative Researcher position at Susquehanna. I am currently a PhD candidate in Economics at Johns Hopkins University. My training combines mathematical reasoning, computational modeling, and data driven empirical analysis. I am motivated by problems that require precise logic, careful implementation, and constant validation using large and noisy datasets. Susquehanna’s emphasis on rigorous research, collaboration with traders and technologists, and long-term strategy reliability aligns closely with how I approach quantitative work.

My research agenda is at the intersection of wealth inequality and macroeconomics. In my job market paper, I build and estimate a structural heterogeneous agent model with returns that differ across households. The core contribution is to show how heterogeneity in returns and income risk jointly shape the wealth distribution and consumption behavior. Relative to closely related work, my framework performs better at matching the lower and middle portions of the wealth distribution. This difference matters quantitatively. Models that focus primarily on returns heterogeneity can fit the top of the wealth distribution closely but tend to generate too much wealth among the bottom half of households. In my framework, income uncertainty generates a precautionary saving motive that produces genuinely low wealth households. This leads to a realistic distribution of marginal propensities to consume and aggregate consumption responses that align with empirical estimates.

A key feature of my research is the heavy use of computation to solve and simulate models that cannot be handled analytically. I write code to solve dynamic programming problems, simulate outcomes under stochastic environments, and evaluate model performance across specifications. I work with large datasets, design backtesting exercises, and focus on robustness and interpretability. This work requires disciplined reasoning, attention to statistical assumptions, and careful validation, skills that map directly to alpha research and strategy evaluation.

I am particularly interested in Susquehanna because of its research driven trading culture and its close integration of researchers, traders, and technologists. The opportunity to focus on model development, robustness, and long-term reliability while collaborating on execution and deployment is especially appealing. I value environments where ideas are stress tested rigorously and where performance provides clear feedback.

My economics training has provided a strong foundation in probability, optimization, and statistical inference. I am comfortable translating theory into working code and analyzing large scale datasets using Python and related tools. I value clear communication and collaborative problem solving, both of which are essential in fast paced and highly interactive research environments.

I would welcome the opportunity to contribute my skills and perspective to Susquehanna. Thank you for your time and consideration.

\closing{Sincerely,}

\end{letter}
\end{document}
